\section{Cylinder experiments and extremal configurations}
\label{sec:cylinder}

This section defines the two cylinder models and constructs the configurations
used in the cylinder estimates.  A state is measured by its squared Hilbert
norm.  Time subdivision copies that state into its descendant intervals;
it does not divide its energy among them.  This distinction explains the
growth factor in the normalization below.

\subsection{States, tests, and assignments}

Fix $0<\eta<1$ and $s=2^{-n}$.  A cut of length $\ell=2^{-k}$, with
$0\le k\le n$, consists of the dyadic subintervals of $[0,1)$ of that
length.  The final cut has length $s$ and contains exactly $s^{-1}$
intervals.  We measure an interval
at its left endpoint $T$.  There are finitely many operator cuts, in
nonincreasing order of length; repeated cuts are permitted.  Their number
is bounded on each polynomial setup.  Every child receives the freely
evolved state of its parent, with the operators specified on that branch.

The exact base index is $F=(X,B,Y)$, with
\begin{equation}\label{cyl:base-flow}
 B_T=B+TX,\qquad Y_T=Y+TX^2.
\end{equation}
Additional exact indices range over standard Borel spaces; arbitrary
separable Hilbert multiplicities are allowed. All integrals and squared
norms include these indices. One may equivalently use a measurable
direct integral of separable Hilbert spaces over such an index space.

\begin{definition}[Cylinder states]\label{cyl:def-states}
In the \emph{classical cylinder}, $(N,D)$ is another exact index,
$D_T=D+TN$, and free evolution leaves the state unchanged in these moving
coordinates.  In the \emph{hybrid cylinder}, $H>0$ is fixed, there is an
exact index $\sigma\in[1,2]$, and a fiber state belongs to
$L^2(\R_D;\mathcal H)$.  Its free propagator $U_t^{H,\sigma}$ is
\[
 \widehat{U_t^{H,\sigma}g}(p)
       =e^{-itHp^2/(2\sigma)}\widehat g(p).
\]
At length $\ell$, set
\[
 d_\ell=(H/\ell)^{1/2},\qquad q_\ell=(H\ell)^{1/2},
 \qquad
 \phi_{D,N}^{\ell,\sigma}(z)
  =(\pi q_\ell^2)^{-1/4}
    e^{-(z-D)^2/(2q_\ell^2)+i\sigma Nz/H}.
\]
The analysis operator is
\[
 (V_\ell g)(D,N)=\ip{g}{\phi_{D,N}^{\ell,\sigma}},
 \qquad d\nu_\sigma(D,N)=\frac{\sigma}{2\pi H}\,dD\,dN.
\]
For the classical model, $V_\ell=\Id$ and $d_\ell=0$.
An admissible operator at a cut is $V_\ell^*MV_\ell$, where $M$ is a
measurable pointwise contraction on the multiplicity space and preserves
all exact indices.  The row measure $\mu_I$ is the squared norm of the
analyzed state, possibly restricted by row assignments.
\end{definition}

With the Fourier normalization of Section~\ref{sec:introduction},
Plancherel in $N$ gives
\[
 \int_\R\abs{\ip{g}{\phi_{D,N}^{\ell,\sigma}}}^2
                   \frac{\sigma\,dN}{2\pi H}
 =\int_\R\norm{g(z)}_{\mathcal H}^2
             (\pi q_\ell^2)^{-1/2}e^{-(z-D)^2/q_\ell^2}\,dz.
\]
Integration in $D$ proves $V_\ell^*V_\ell=\Id$.  In particular, every
admissible operator is a contraction.  The packet parameter $D$ in the
hybrid model is an observation coordinate, not an exact classical index.

\begin{definition}[Cylinder tests]\label{cyl:def-test}
A test $Q$ at $(T,\ell)$ has widths $r>0$, $f\ge d_\ell$, with $f>0$.
For a hybrid row observed at the cut $T$, $D_T$ in the following
inequalities denotes its instantaneous analyzed center $D$, not an
exact classical trajectory. The test is specified by
\begin{gather}
 |X-x|\le r,\qquad |B_T-b|\le\ell r,\qquad
 |Y_T-y-2x(B_T-b)|\le\ell r^2,\label{cyl:base-test}\\
 |N-c-\alpha X-\beta(T)X^2-K(Y_T-2XB_T)|\le f,
 \label{cyl:velocity-test}\\
 |D_T-d-\alpha B_T-\beta(T)Y_T+KB_T^2|\le\ell f.
 \label{cyl:position-test}
\end{gather}
All centers and coefficients are real.  Extend the coefficients by
\[
 \beta(t)=\beta(T)+(t-T)K,\qquad d(t)=d+(t-T)c.
\]
The test weight is $w(Q)=r^{3+\eta}f^{1-\eta}$.  A row measure is
$\kappa$-regular if $\mu(Q)\le\kappa w(Q)$ for every admissible test.
\end{definition}

If the expressions subtracted from $N,D_T$ in
Equations~\eqref{cyl:velocity-test}--\eqref{cyl:position-test} are
$A(F)$ and $S_t(F)$, respectively, direct differentiation gives
$\partial_t A=0$ and $\partial_t S_t=A$.  Also
\begin{equation}\label{cyl:derivative}
 P(F)=\alpha+2\beta(t)X-2KB_t
\end{equation}
is independent of $t$.

An assignment to the final cut is a family of nonnegative row measures
$\mu_\lambda\le\mu_I$, each supported on a test $Q_\lambda$ at its own
interval, such that $\sum_{\lambda\text{ at }I}\mu_\lambda\le\mu_I$.
Fractional assignments are allowed.  Put $m_\lambda=\mu_\lambda(Q_\lambda)$
and $w_\lambda=w(Q_\lambda)$.  If $e>0$ is the root energy and the root
row measure is $\kappa$-regular, define
\begin{equation}\label{cyl:functional}
 \mathcal B_s
  =\frac{s^{-1/2}}{e\sqrt\kappa}
           \sum_\lambda\frac{m_\lambda^{3/2}}{\sqrt{w_\lambda}}.
\end{equation}
The sum includes all final intervals.  Bounded enlargements of the defining
inequalities are interchangeable with a bounded change of the weight:
enlarge $r,f$ by fixed constants or cover by finitely many translates.

\subsection{Cutoffs and localization}

\begin{definition}[Polynomial cylinder setup]\label{cyl:def-setup}
The quantities $e/\kappa,\kappa/e$ and, in the hybrid case, $H,H^{-1}$
are bounded by fixed powers of $s^{-1}$.  The initial base support lies in
a box of fixed-power size.  The initial normal phase has negligible mass
relative to $e$ outside such a box.  The marginal base density and the
joint density in $(X,B,Y,N,D)$ are bounded by $e$ times a fixed power of
$s^{-1}$, with respect to Lebesgue measure.  Negligibility is uniform
within each family to which an exponent bound is applied.  All these
powers, and the bound on the number of operators on a branch, are fixed
before $s\to0$.
\end{definition}

The exponent closure may use successively larger fixed powers, in the
ordered sense specified in the framework.  It never replaces the inner
scale limit by isolated scales with increasing cutoff powers.

\begin{lemma}[Symmetries and Gaussian localization]\label{cyl:localization}
The models are invariant under translations, base boosts, dilation in
base units $(r,r,r^2)$, normal dilation, affine time normalization, and
subtraction of one common normal shear of
Equations~\eqref{cyl:velocity-test}--\eqref{cyl:position-test}.  All
operators and assignments must be transformed with the state.

For $\ell'\le\ell$, $|t|\le C\ell$, and fixed $\eps>0$, the hybrid
frame transfer is negligible in operator norm outside
\begin{equation}\label{cyl:correspondence}
 |N'-N|\le s^{-\eps}d_{\ell'},\qquad
 |D'-D-tN|\le s^{-\eps}q_\ell.
\end{equation}
This remains true along any fixed finite stack of masks and after summing
over a polynomial number of intervals.  At a repeated cut, synthesis and
reanalysis of a restricted $\kappa$-regular row measure produce a
$C_\eta\kappa$-regular measure.  The same assertion holds for truncated
families of tests whenever the normal enlargements in the proof belong
to that family or are covered by translates belonging to it.
\end{lemma}

\begin{proof}
A base boost replaces $X$ by $X-x_0$ and replaces the last base position
by its appropriate linear tilt in $B_t$; Equation~\eqref{cyl:base-flow}
then has the same form.  A common normal shear subtracts
$a(F)$ from velocity and $b(F)+ta(F)$ from position.  On each hybrid
fiber it is a Weyl translation, and the free propagator intertwines the
two translations up to a scalar phase.  The scalar phase can be included
in the conjugated mask.  If a time interval of length $L$ is normalized
to unit length, positions are divided by $L$, velocities are unchanged,
and $H$ becomes $H/L$.  Dividing normal coordinates and velocities by
$F_0$ changes $H$ to $H/F_0^2$.  Dividing base coordinates by
$(r_0,r_0,r_0^2)$ and normal coordinates by $F_0$ divides all weights
by $A=r_0^{3+\eta}F_0^{1-\eta}$ and multiplies the regularity constant
by $A$.  Energy is preserved by the corresponding unitary coordinate
changes.  Thus Equation~\eqref{cyl:functional} is unchanged.

The scalar transfer kernel is
$\ip{U_t^{H,\sigma}\phi_{D,N}^{\ell,\sigma}}
          {\phi_{D',N'}^{\ell',\sigma}}$.
Free propagation changes the first packet's position width to at most
$Cq_\ell$, since $|t|H/q_\ell\le Cq_\ell$, and leaves its Fourier
width comparable to $q_\ell^{-1}$.  Integrating two Gaussians, in physical
space for the position difference and in Fourier space for the momentum
difference, gives an upper bound of the form
\[
 C_s\exp\left[-c\left(
       \frac{|D'-D-tN|^2}{q_\ell^2}
       +\frac{|N'-N|^2}{d_{\ell'}^2}\right)\right],
\]
where $C_s$ and the Schur integration volumes have only fixed-power
costs on the setup.  Each of the two separate Gaussian bounds implies
the displayed joint bound after reducing $c$.  Schur's test therefore
makes the discarded kernel smaller than every power of $s$.  A finite
product and a polynomial sum preserve this conclusion.  If no chain of
the correspondences connects two row sets, the operator between those
sets is negligible; this statement does not require a positive-kernel
description of its main part.

At a repeated cut the absolute kernel is exactly
\[
 \exp\left(-\frac{(D-D')^2}{4q_\ell^2}
           -\frac{\sigma^2(N-N')^2}{4d_\ell^2}\right).
\]
Its Schur norms are bounded independently of $H,\ell$.  Cauchy--Schwarz
against this kernel bounds output energy on a test by a Gaussian-weighted
sum of input energies on its normal enlargements.  Since $f\ge d_\ell$
and $\ell f\ge q_\ell$, the enlargement in the $j$th Gaussian shell
is obtained by multiplying the normal width by $C(1+j)$.  The resulting
regularity bound is
$C\kappa w(Q)\sum_{j\ge0}e^{-cj^2}(1+j)^{1-\eta}$.
Exact base indices are never mixed.  This proves the last assertion.
\end{proof}

\begin{lemma}[Polynomial test reduction]\label{cyl:polynomial-tests}
At every stage of a polynomial setup, the base and normal cutoff and
density conditions hold with enlarged fixed powers.  To any prescribed
polynomial regularity floor or error tolerance, it suffices to test
widths, reciprocal widths, and absolute test parameters bounded by
fixed powers.  Discarding output tests whose mass-to-weight ratio is
below $\rho$ costs at most $\sqrt\rho$ times total assigned mass in
the cubic readout, independently of the number of tests.
\end{lemma}

\begin{proof}
Base indices move exactly and masks are contractions on each fiber.
In the hybrid model, the pointwise packet bound gives a joint phase
density at most $C H^{-1}$ times the marginal base energy density.
Lemma~\ref{cyl:localization} propagates normal tails.  The density and
base-support bounds give, uniformly in shear parameters,
\begin{equation}\label{cyl:volume-bound}
 \mu(Q)\le e s^{-C}\min(r^4,1)\min(f^2,1).
\end{equation}
For $f\le1$ use joint density, and for $f>1$ use marginal base density.
Dividing by $r^{3+\eta}f^{1-\eta}$ shows that either sufficiently small
width gives an arbitrarily small prescribed power: the factors are
$r^{1-\eta}$ and $f^{1+\eta}$, respectively.  If neither width is small
and one is sufficiently large, the total-energy bound suffices.

We record an elementary sublevel estimate to handle coefficients.  For a
nonzero real polynomial $Q$ of degree at most two in $d\le3$ variables, on a
fixed cube,
\begin{equation}\label{cyl:sublevel}
 \vol\{\abs Q\le\rho\norm{Q}_{\mathrm{coeff}}\}
       \le C_d\rho^{2^{-d}},\qquad 0<\rho<1.
\end{equation}
Here the coefficient norm is the maximum absolute coefficient.  In one
variable, three points in a set of measure $m$, separated by at least
$cm$, and Lagrange interpolation show
$\norm Q_{\mathrm{coeff}}\le C\rho\norm Q_{\mathrm{coeff}}/m^2$.
For the induction, write $Q$ as a polynomial in the last variable and
choose a coefficient polynomial with coefficient norm equal to
$\norm Q_{\mathrm{coeff}}$.  Outside the set where this coefficient is
smaller than $\rho^{1/2}\norm Q_{\mathrm{coeff}}$, the one-dimensional
bound is $C\rho^{1/4}$; the exceptional base set has measure at most
$C\rho^{2^{-d}}$ by induction.  This proves
Equation~\eqref{cyl:sublevel}, with harmless changes of constants.

For widths in a fixed polynomial range, discard negligible normal tails
and clip base centers to a fixed-power box.  The two normal shear
polynomials have coefficient norm comparable, by a fixed-power change
of coordinates, to the maximum of $|c|,|d|,|\alpha|,|\beta(T)|,|K|$.
If that maximum is sufficiently large, at least one of their simultaneous
sublevel conditions has arbitrarily small prescribed polynomial volume
by Equation~\eqref{cyl:sublevel}.  The density bound handles its mass.
Finally,
$m_\lambda^{3/2}/\sqrt{w_\lambda}
 =m_\lambda\sqrt{m_\lambda/w_\lambda}$, so small ratios are summed
against assigned mass, not against label cardinality.
\end{proof}

We shall also use the following convention for very small pieces.  An
experiment with input energy $e'$ and prescribed regularity $\kappa'$
satisfies the direct estimate
\begin{equation}\label{cyl:tiny-energy}
 \mathcal B_s\le s^{-3/2}
       \max_\lambda\sqrt{e'/(\kappa'w_\lambda)}.
\end{equation}
Indeed $m_\lambda\le e'$ and total assigned mass is at most $s^{-1}e'$.
For uniform application to a family of pieces one uses a fixed threshold,
not a scale-dependent notion of smallness.  Let the old reference energy
be $e_0$, suppose the summed piece energies have a polynomial bound
relative to $e_0$, and let $w_{\min}$ be a common polynomial lower bound
for relevant output weights.  Pieces with $e_G\le s^A e_0$ have total
cubic readout at most
\[
 s^{-1}w_{\min}^{-1/2}\sum_G e_G^{3/2}
 \le s^{-1}w_{\min}^{-1/2}(s^A e_0)^{1/2}\sum_Ge_G.
\]
Taking the fixed $A$ sufficiently large makes this smaller than any
specified polynomial comparison scale.  Every remaining piece has
uniform cutoff and negligible-tail bounds relative to its own energy,
after enlarging fixed powers.  Prescribed regularity constants in the
applications are also polynomial relative to $e_0$, so both energy to
regularity ratios are polynomial.  The uniform model exponent from
Lemma~\ref{lem:fw-uniform} therefore applies to this remaining family.

\subsection{The extremal exponent and ordinary splitting}
\label{cyl:extremal-exponent}

For each cylinder model, take the supremum of the upper
$\log_{1/s}$-exponents of Equation~\eqref{cyl:functional} over its
polynomial setups. Denote these exponents by $\Gamma_{\mathrm{cl}}$
and $\Gamma_{\mathrm{hyb}}$. The uniform interpretation on each fixed
setup, and all closure extrema below, use the ordered conventions of
Section~\ref{sec:framework}.

\begin{theorem}[Cylinder exponent bound]\label{thm:cylinder}
For the classical and hybrid cylinder experiments with the
regularity weight parameter $0<\eta<1$,
\[
 \Gamma_{\mathrm{cl}}\le\eta/2,
 \qquad \Gamma_{\mathrm{hyb}}\le\eta/2.
\]
\end{theorem}

The proof occupies the rest of this section and
Sections~\ref{sec:low}--\ref{act:section}. We first bound the growth
exponent, then study experiments whose cubic readout approaches it.
Any ordinary regularity split that preserves this extremal growth forces its
energy and regularity bounds to be sharp at the level of exponents.
We then choose the retained tests with compatible shears and describe
how their widths vary with depth. These properties supply the geometric
input for the entropy argument.

Fix either model and write $\Gamma$ for its exponent. We begin by proving
that it is finite.

\begin{lemma}[One-time bound]\label{cyl:one-time}
For a polynomial test at a cut of length $\ell$, an experiment with
$\kappa$-regular input satisfies
\begin{equation}\label{cyl:one-time-bound}
 \mu_I(Q)\le s^{-o(1)}\kappa w(Q)\ell^{1+\eta/2}.
\end{equation}
Consequently $\Gamma$ is finite.
\end{lemma}

\begin{proof}
Subtract the test's shear.  Backward packet correspondences place its
normal predecessors in root conditions of width $s^{-o(1)}f$, since
every intervening $d_{\ell'}$ is at most $d_\ell\le f$.  Split its
$X$-range into $O(\ell^{-1/2})$ intervals of width $r\sqrt\ell$.
For each such interval, backflow of the two base position conditions is
covered by a bounded number of root tests of base width $r\sqrt\ell$:
the horizontal error is $O(r\sqrt\ell)$ and the curvature error in the
tilted coordinate is $O(\ell r^2)$.  Summing their weights gives
\[
 O(\ell^{-1/2})(r\sqrt\ell)^{3+\eta}f^{1-\eta}
       =O(1)w(Q)\ell^{1+\eta/2}.
\]
Contraction on this predecessor region and
Lemma~\ref{cyl:localization} prove the claim; polynomial test bounds
absorb the negligible errors.  Total output mass is at most $s^{-1}e$,
so Equation~\eqref{cyl:one-time-bound} also gives
$\mathcal B_s\le s^{-1+\eta/4-o(1)}$.
\end{proof}

Suppose henceforth that $\Gamma>0$.  Dyadic binning of widths and
mass-to-weight ratios produces extremizing families with common widths
and
\begin{equation}\label{cyl:pd}
 \frac{m_\lambda}{\kappa w_\lambda}=s^{p+o(1)},\qquad
 \sum_\lambda m_\lambda=e s^{-d+o(1)},\qquad
 \Gamma=d+\frac{1-p}{2},\quad p\ge1,\quad0<d\le1.
\end{equation}
Thus $p$ measures the decay of each assigned mass relative to its
regularity bound, whereas $d$ measures the total mass read across the
copied time intervals. The identity for $\Gamma$ follows by substituting
the two mass relations into Equation~\eqref{cyl:functional}.
Negligible ratio classes are first removed by
Lemma~\ref{cyl:polynomial-tests}; there are subpower many remaining bins.
Choose the smallest possible $p=p_0$ in the closure of these extremal
data. This selects the largest normalized assigned mass-to-weight ratios
compatible with extremal growth. In particular $p<3$. The output label
family then has
\begin{equation}\label{cyl:terminal-count}
 \sum_{\lambda\ \mathrm{allowed}}w_\lambda
       \le(e/\kappa)s^{-(p+d)+o(1)}.
\end{equation}

\begin{lemma}[Regularity peeling]\label{cyl:peeling}
At a fixed cut, a state can be decomposed coherently into subpower many
row projections followed by synthesis.  Apart from a remainder at any
prescribed polynomial floor, each level has a number $b>0$ with the
following properties.  Its raw row measure is $O(b)$-regular; it has
disjoint assignments to tests of masses between $b w$ and $C b w$; if
$z$ is its total raw energy, their total weight is at most $z/b$.
The synthesized energies sum to at most $z$, and each synthesized state
is $s^{-o(1)}O(b)$-regular.  Levels may be subdivided by subpower many
statistics without altering these upper bounds.
\end{lemma}

\begin{proof}
Start above the fixed-power upper bound furnished by
Lemma~\ref{cyl:polynomial-tests}, and descend through dyadic thresholds.
At threshold $b$, remove the entire remaining part of a test whenever its
mass exceeds $bw$.  Its mass is at most $2bw$, by regularity of the
preceding remainder.  A polynomial lower bound on the allowed weights
makes the removal procedure finite by total mass.  The union removed at
this threshold is dominated by that preceding remainder and is therefore
$2b$-regular.  The unremoved measure is $b$-regular.  Continue to the
prescribed floor.  The removed row sets are disjoint and their projections
sum to the identity.  Synthesis is a contraction, and
Lemma~\ref{cyl:localization} gives post-synthesis regularity.

For a fixed terminal assignment, its cubic readout to the power $1/3$
is a weighted $\ell^3$ norm of Hilbert norms.  Hence coherent decomposition
costs at most a fixed power of the number of levels by the triangle
inequality.  That number is subpower, and at least one level preserves
any extremal output exponent.
\end{proof}

\begin{proposition}[Saturation of a preserving split]\label{cyl:split-saturation}
Retain the terminal family in Equation~\eqref{cyl:terminal-count}.
After any finite sequence of inserted operators or splittings that
preserves its extremal exponent, every dominant homogeneous terminal
class again has the exponents $p,d$ in Equation~\eqref{cyl:pd}.
Suppose a non-remainder level supplied by Lemma~\ref{cyl:peeling} at
depth $a\in(0,1)$ preserves this exponent. Let $b$ be its peeling
threshold and $z$ its total raw energy, summed over the intervals of
that cut. Then
\begin{equation}\label{cyl:split-benchmarks}
 b=\kappa s^{pa+o(1)},\qquad
 z=e s^{-da+o(1)}.
\end{equation}
The summed post-synthesis energy has the same exponent as $z$.
The conclusion holds for a further sublevel that preserves the output.
\end{proposition}

\begin{proof}
A new terminal class has $p'\ge p$ by minimality and
$d'=\Gamma+(p'-1)/2$.  Its assigned mass and the fixed weight count imply
$p'+d'\le p+d$.  Both inequalities force $p'=p$ and $d'=d$.

For a non-remainder peeling level at $\ell=s^{a+o(1)}$, apply the prefix
exponent to its raw heavy-test assignments.  Lemma~\ref{cyl:peeling}
gives
\begin{equation}\label{cyl:prefix-split}
 \ell^{-1/2}z\sqrt{b/\kappa}
       \le \ell^{-\Gamma-o(1)}e.
\end{equation}
The tail exponent, applied separately to each interval and summed with
its own post-energy, gives cubic readout at most
\[
 (s/\ell)^{1/2-\Gamma-o(1)}z\sqrt b.
\]
Combining this with Equation~\eqref{cyl:prefix-split} recovers the full
extremal upper exponent.  Since the chosen level preserves that exponent,
both estimates, and the synthesis energy bound, must be saturated in
exponent.  The floor remainder is negligible: its summed input energy
is at most $\ell^{-1}e$, and a sufficiently small regularity floor makes
the tail readout smaller by a fixed power.

Write $b/\kappa=s^{P+o(1)}$.  The saturated prefix has mass-to-weight
exponent $P/a$, so minimality implies $P/a\ge p$.  For the tail first
select a homogeneous dominant terminal mass class, then an interval
whose normalized readout is near maximal.  That tail realizes $\Gamma$
in the closure and has mass-to-weight exponent $(p-P)/(1-a)$; hence
$(p-P)/(1-a)\ge p$.  The one-time bound bounds $P$ on the fixed gap, so
these are legitimate finite closure data.  Thus $P=pa$, and saturation
in Equation~\eqref{cyl:prefix-split} gives $z=e s^{-da+o(1)}$.

For a sublevel, use the unpartitioned level in the prefix argument.
It bounds the sublevel's raw energy from above.  The preserving tail
bounds its post-energy from below by the same exponent.  At the root,
the given input regularity and energy replace the prefix estimate.
At the last cut a repeated-cut regularity estimate replaces the tail.
\end{proof}

For later use put
\begin{equation}\label{cyl:benchmarks}
 \ell_i=s^i,\qquad E_i=e s^{-di},\qquad\kappa_i=\kappa s^{pi}.
\end{equation}
Depths and widths are dyadically rounded.  All exponent upper and lower
bounds below permit $s^{\mp o(1)}$ factors.  On a fixed gap $h>0$, these
are also vanishing losses in $\log_{s^h}$ units.

\subsection{Matched gates}

Ordinary splitting controls energy and regularity at successive cuts.
We now also align the shear frames of the retained tests. This compatibility
will allow a later phase name to be compared with earlier names along a
retained path.

A \emph{documented gate} is a row projection followed by synthesis,
whose retained rows are partitioned among tests.  A retained row belongs
to its assigned test; the associated test is its documented label.
The last gate is simply the final assignment and need not be synthesized.
Statements about reachability refer to chains of
Equation~\eqref{cyl:correspondence} through the fixed retained
mask/test row sets. They do not impose the original state's
positive-norm support at intermediate cuts, which need not persist
after an input restriction. They assert geometric incidence, not a
joint probability law for the input and output of a quantum operator.

\begin{proposition}[Matched-chain construction]\label{cyl:matched-chain}
On any fixed finite depth mesh $0\le i_0<\cdots<i_b=1$, an extremizing
family can be modified by finitely many masks, preserving its terminal
exponent on the original label family, so that the following hold.
\begin{enumerate}[label=\textup{(\roman*)}]
\item At depth $i$, the documented tests have common widths $r_i,f_i$
and total weight at most $s^{-o(1)}E_i/\kappa_i$.
\item Ordinary regularization may be imposed immediately before each
gate, in forward order, giving summed energy at most $s^{-o(1)}E_i$
and regularity at most $s^{-o(1)}\kappa_i$ in each interval.
\item The total mass read on the documented rows is at least
$s^{o(1)}E_i$.  This lower bound persists after any further restrictions
that preserve the terminal exponent and retain these gates in series.
\item If $j$ is the next mesh depth and $\delta=s^{j-i}$, then, for
$L_*=L_*(\eta)$ independent of the mesh,
\[
 \delta^{L_*}\lesssim r_i/r_j\le1,\qquad
 \delta^{L_*}\lesssim f_i/f_j\lesssim\delta^{-L_*}.
\]
On every connecting packet chain between the two documented labels,
\begin{equation}\tag{M}\label{eq:M}
 \begin{split}
 |\beta_i(T_i)-\beta_j(T_i)|+\ell_i|K_i-K_j|
       &\lesssim(f_j/r_j^2)\delta^{-L_*},\\
 |P_i(F)-P_j(F)|&\lesssim(f_j/r_j)\delta^{-L_*}.
 \end{split}
\end{equation}
These comparisons survive insertion of any fixed finite intervening
operator stack, with a sufficiently small correspondence truncation
exponent relative to the mesh gaps.
\end{enumerate}
\end{proposition}

We first prove the geometric estimate used in its construction.

\begin{lemma}[Two-shear intersection]\label{cyl:intersection}
Normalize one starting interval and the later test widths to
$\ell_i=r_j=f_j=1$, and put $\delta=\ell_j/\ell_i$.
Let $Q^-,\widetilde Q^-$ be two projected tests, with bounded base
footprints and normal widths $O(\delta^{-1})$.  At the first base center
let $M$ be the maximum of $1$ and the differences of $P,\beta(T_i),K$.
For a $\kappa_i$-regular starting measure,
\begin{equation}\tag{I}\label{eq:I}
 \mu(Q^-\cap\widetilde Q^-)
       \le s^{-o(1)}\delta^{-C}\kappa_i w_j/M.
\end{equation}
If every retained descendant label has mass at least
$\delta^D\kappa_iw_j$, then a fixed projected label has at most
$s^{-o(1)}\delta^{-C_D}M_0^{1-\eta}$ intersecting partners in the
shell $M\asymp M_0$.
\end{lemma}

\begin{proof}
Subtract the first shear and center its base chart.  The intersection
imposes, on bounded $(X,B,Y)$,
\[
 |\alpha X+\beta X^2+k(Y-2XB)+c'|\lesssim\delta^{-1},\qquad
 |\alpha B+\beta Y-kB^2+d'|\lesssim\delta^{-1},
\]
where $\max(1,|\alpha|,|\beta|,|k|)\asymp M$.
If $\max(|\beta|,|k|)\gtrsim M$, multiplying the first expression by
$\beta$ and subtracting $k$ times the second gives
\[
 |t^2+\alpha t+\mathrm{const}|\lesssim M\delta^{-1},
 \qquad t=\beta X-kB.
\]
On a dyadic shell $R\asymp\max(\sqrt M,|2t+\alpha|)$, its admissible
$t$-set lies in at most two intervals of length $O(M\delta^{-1}/R)$.
In the lowest shell the same conclusion follows from proximity to the
quadratic vertex.  Since $\max(|\beta|,|k|)\gtrsim M$, the horizontal
set is covered by $O(\delta^{-1}R)$ squares of side $1/R$.

In such a square use the coordinate $Y-2x_{\mathrm{cell}}B$.
After the vertical terms are removed, the horizontal derivatives of the
two expressions are $O(R+M/R)$, so their oscillation across the square
is $O(1)$ because $R\ge\sqrt M$.  At least one vertical coefficient has
size comparable to $M$.  Therefore
$O(\delta^{-1}R^2/M)$ vertical intervals of length $R^{-2}$ suffice.
The intersection is covered by
$O(\delta^{-2}R^3/M)$ Heisenberg base cells of width $O(1/R)$, together
with the first normal conditions enlarged by $O(\delta^{-1})$.
Regularity bounds their total mass by
\[
 C\delta^{-C}\kappa_i w_j
       \frac{R^3}{M}R^{-3-\eta}
       \le C\delta^{-C}\kappa_iw_j/M.
\]
There are only logarithmically many shells on a polynomial setup.

Otherwise $|\alpha|\asymp M$ and $|\beta|,|k|$ are sufficiently small
relative to $M$ on the bounded base box.  For fixed $Y$, the second
expression has a $B$ derivative comparable to $M$, and the first then
has an $X$ derivative comparable to $M$.  The volume is
$O(\delta^{-2}M^{-2})$.  Its Heisenberg neighborhood of width $1/M$
obeys the same estimate, since the expressions change only by a constant
there.  It can consequently be covered by
$O(\delta^{-2}M^2)$ cells of weight $O(M^{-3-\eta})$, proving the same
bound.  Bounded $M$ needs no subdivision.

For the partner count, every intersecting partner in the shell has its
whole projected footprint in a bounded enlargement of the first base
box, and in normal width $O(\delta^{-1}M_0)$ about the first shear.
This follows by extending the difference polynomials from an intersection
point across bounded base coordinates.  At each descendant time,
contraction and starting regularity bound the total mass of these
assigned rows by
$s^{-o(1)}\kappa_i w_j O(\delta^{-1}M_0)^{1-\eta}$.
Divide by the assumed mass per label and sum over at most $O(\delta^{-1})$
descendant times.
\end{proof}

\begin{proof}[Proof of Proposition~\ref{cyl:matched-chain}]
\textbf{Read mass at a gate.}
First observe that the asserted mass lower bound follows whenever the
gate and weight count have been constructed.  Insert an ordinary split
immediately after the gate in a thought experiment preserving the
terminal exponent.  Proposition~\ref{cyl:split-saturation} requires
summed post-energy $s^{o(1)}E_i$, and contraction bounds it by the raw
gate mass.  At the final cut use the fixed terminal count and minimality
of $p$.  At the root, repeated-cut regularity is $s^{-o(1)}\kappa$,
so reducing post-energy by a fixed power would lose the terminal exponent
by the full tail bound.

\paragraph{Removing incompatible rows.}
Construct the gates backwards.  Suppose the $j$-gate is fixed and set
$a=j-h$, $\delta=s^h$.  At $a>0$ make an ordinary preserving split;
at $a=0$ use the root state.  The resulting states $g_a$ have total
energy at most $s^{-o(1)}E_a$ and regularity at most
$s^{-o(1)}\kappa_a$.  At $j$ their row measure is at worst
$s^{-o(1)}\kappa_a$-regular by Lemma~\ref{cyl:one-time}.
Delete labels whose newly read mass is less than
$\delta^D\kappa_aw_j$.  Their total mass is at most
\[
 s^{-o(1)}\delta^D\kappa_a E_j/\kappa_j.
\]
After synthesis they remain $s^{-o(1)}\kappa_a$-regular.  Their tail
readout, relative to the required scale $E_j\sqrt{\kappa_j}$ at $j$,
therefore loses at least $\delta^{D-3p/2-o(1)}$.  Since $p<3$, an
absolute $D$ makes the deletion harmless.  At $j=1$, use the output
ratio bound directly.

Every retained descendant test projects to a starting test $Q^-_\lambda$
with base width $O(r_j)$ in $\ell_a$ position units and normal width
$O(\delta^{-1}f_j)$.  Base backflow follows from the quadratic identity;
normal backflow follows from Lemma~\ref{cyl:localization}, taking its
truncation exponent smaller than $h/10$.  Each such projected test
contains every predecessor on the allowed correspondence chains.
All label counts here are polynomial by their weight count and the
polynomial lower bound for $w_j$.

Call a starting row ambiguous if it lies in two projected tests with
$M>\delta^{-D_0}$.  Lemma~\ref{cyl:intersection}, its partner count,
and the $j$-gate weight count bound its total mass by
\[
 s^{-o(1)}E_j(\kappa_a/\kappa_j)
                  \delta^{-C}\delta^{\eta D_0}.
\]
Choose $D_0=D_0(\eta)$ large enough that this is smaller than $E_a$ by
a fixed power of $\delta$.  Restricted rows retain regularity after
synthesis, so their tail is negligible by the model exponent.  Remove
them and also rows belonging to no projected retained test.  The latter
have negligible transfer to the $j$-gate.  Both removals use one row
projection followed by synthesis, not successive assumptions about
disjointness of synthesized states.

\paragraph{Shear charts.}
Partition the remaining rows into charts $G$.  First record a lattice
Heisenberg base cell of width $r_j$.  Choose one attached test in a fixed
order and round its five shear data at that cell center: $\beta(T_a)$,
$K$, $P$, and its two normal values.  Round in normal units
\[
 f^*=\delta^{-D_*}f_j,\qquad D_*>D_0+2.
\]
The five data determine the polynomial by a triangular change of
coordinates.  Nonambiguity makes every attached test's data lie within
a bounded number of these rounded values.  Thus each $j$-label meets
only boundedly many charts.  In each chart, after subtracting the common
rounded shear, base and normal support are bounded in
$(r_j,f^*,\ell_a)$ units, and every attached $j$-test has bounded
coefficients in those units.

\paragraph{Full chart regularity.}
The charts now have compatible geometry. To apply the model exponent to
their synthesized states, we also need full regularity. The next peel gives
regularity on a bounded range of tests; the subsequent reanalysis,
smoothing, and cutoff estimates extend it to the full test family.

Within each chart perform a truncated peel: normalized widths belong to
$[\delta^{D_1},1]$, with the Planck floor retained, and normalized shear
coefficients have absolute value at most $\delta^{-D_1}$.  Centers are
unrestricted.  The physical weight multiplier is
$w^*=r_j^{3+\eta}(f^*)^{1-\eta}$.  Keep one density and width class
coherently across all charts; there are subpower many classes, so some
class preserves the output.  Let $b_*$ be its density bound, or its
floor if it is the remainder.  Then $b_*\le s^{-o(1)}\kappa_a$,
the raw mass in each chart is $O(b_*w^*)$, and the synthesized chart
energies satisfy $\sum_G e_G\le s^{-o(1)}E_a$.

We justify applying the full model exponent with regularity $b_*$ to
each chart.  Same-cut Gaussian reanalysis preserves its truncated
regularity: normal shell enlargements can be covered by translates at
unchanged widths and coefficients.  Smooth the input and the entire
subexperiment by independent base boosts, base position translations,
and the two normal translations, each of size $\delta^l$, where
$l>D_*+10$.  Retain the smoothing parameter as an exact auxiliary index.
These are exact symmetries.  The smallest attached output accuracy is
$\delta^{1+D_*}$, so assigned masses persist on bounded enlarged tests.
In normalized units the smoothed base and joint densities are at most
\[
 C\delta^{-5l} b_* w^*.
\]
This follows by averaging translations, whose coordinate Jacobians are
one, and using the raw chart mass bound.  Base support is bounded;
normal support outside $s^{-o(1)}$ is negligible by Gaussian tails.

Here are explicit choices showing that the fixed density cost disappears
from the regularity estimate.  Choose $D_2$ with
$(1-\eta)D_2>5l+1$.  Equation~\eqref{cyl:volume-bound} gives
\[
 \frac{\mu(Q)}{b_*w^*w(Q)}
 \le C\delta^{-5l}
 \frac{\min(r^4,1)\min(f^2,1)}{r^{3+\eta}f^{1-\eta}}.
\]
Thus any width below $\delta^{D_2}$ is harmless.  Widths above one are
handled by bounded base covering or, for the normal width, by dropping
the normal conditions and covering the $s^{-o(1)}$ normal support with
unit translates.  For widths in $[\delta^{D_2},1]$, choose $D_3$ so
that $2^{-3}D_3>5l+4D_2+1$.  Equation~\eqref{cyl:sublevel} makes tests
with a coefficient above $\delta^{-D_3}$ harmless, even after dividing
by their minimum weight $\delta^{4D_2}$.  All remaining tests pull back
under each jitter to permitted truncated tests, provided
$D_1>D_2+D_3+10$.  For them the original truncated regularity bound is
used directly and incurs no factor $\delta^{-5l}$.
This proves full regularity $s^{-o(1)}b_*w^*$ in normalized units.
Tiny-energy chart pieces use Equation~\eqref{cyl:tiny-energy}; the others
satisfy the polynomial setup relative to their own energy.

\paragraph{Saturation and induction.}
The model exponent and the bounded chart overlap now bound the coherent
cubic readout at $j$ by
\begin{equation}\label{cyl:chart-upper}
 s^{-o(1)}\delta^{1/2-\Gamma}\sqrt{b_*}\sum_Ge_G.
\end{equation}
Indeed, for a label reached by at most $C$ chart states,
$\norm{\sum_{G=1}^Cv_G}^3\le C^2\sum_G\norm{v_G}^3$.
The already established read-mass lower bound and weight count, by
H\"older, give the lower bound
\[
 s^{o(1)}E_j\sqrt{\kappa_j}
  =s^{o(1)}\delta^{1/2-\Gamma}E_a\sqrt{\kappa_a}.
\]
Consequently $b_*=\kappa_a s^{o(1)}$, and the floor remainder can be
excluded.  Use the heavy tests in the selected truncated level as the
new documented labels.  Their total weight is at most
$s^{-o(1)}E_a/\kappa_a$.  Their truncated widths and coefficients,
together with chart compatibility, give the width bounds and
Equation~\eqref{eq:M}.  The constants depend only on the choices above,
hence only on $\eta$, not on the number of mesh steps.  Any reachable
$j$-test is attached to the actual supporting row of the new gate; no
assertion about its mass in an earlier state is used.

This completes the backward induction.  Finally insert ordinary
regularization before the gates in forward order, preserving the
terminal exponent at every step.  Proposition~\ref{cyl:split-saturation}
gives the desired benchmarks.  At the final cut its prefix minimality
and the repeated-cut regularity bound give the same conclusion; the
root uses its given bounds.  Additional finite row statistics may be
flattened in this same forward order.  Dyadic upper endpoints may be
clipped at $r_j$, so repeated rounding does not create exponential
drift in $r_i/r_j\le1$.  Constants per step disappear before the mesh
limit because every inner mesh is finite.
\end{proof}

\subsection{Extraction and the hybrid Planck floor}

\begin{lemma}[Group extraction]\label{cyl:group-extraction}
Consider a fixed positive depth gap $[i,j]$ in a documented configuration,
with ordinary pre-gate regularity at $i$.  Partition its starting rows
into groups, including their starting intervals, and synthesize each
group separately.  Suppose every used $j$-label is reachable from at
most $s^{-o(1)}$ groups; all other transfers must be negligible by uniform
packet locality.  Then one can extract group subexperiments realizing
$\Gamma,p$ in the closure.  Their reachable output weight satisfies
\[
 \sum_{\lambda\text{ reachable from }G}w_\lambda
   \le \frac{e_G}{\kappa_i}
                 s^{-(p+d)(j-i)-o(1)}.
\]
The same conclusion can be imposed simultaneously at finitely many
intermediate documented depths with the corresponding overlap bound.
\end{lemma}

\begin{proof}
The group energies sum to at most $s^{-o(1)}E_i$, and each group has
regularity $s^{-o(1)}\kappa_i$ after reanalysis.  H\"older, gate count,
and read mass give cubic output at least
$s^{o(1)}E_j\sqrt{\kappa_j}$.  The overlap assumption compares this to
the sum of the group readouts with subpower loss.  Each group readout
is at most
$s^{(1/2-\Gamma)(j-i)-o(1)}e_G\sqrt{\kappa_i}$.
Thus near-extremal normalized groups must exist.

Summing reachable weights over groups costs at most
$s^{-o(1)}E_j/\kappa_j$.  Groups violating the displayed relative count
with an extra factor $s^{-\eps}$ have total energy at most
\[
 s^{(p+d)(j-i)+\eps-o(1)}\kappa_i E_j/\kappa_j
       =s^{\eps-o(1)}E_i.
\]
Their readouts are negligible by the group upper estimate.  Delete them,
then extract a near-maximal remaining group and let $\eps\downarrow0$.
Its terminal homogeneous class has exponent $p$ by the same minimality
and count argument as Proposition~\ref{cyl:split-saturation}.
Finitely many intermediate count deletions are handled in the same way.

Only polynomially many groups are relevant: there are polynomially many
labels and at most subpower many groups reaching each.  Remove globally
unreachable groups as one complementary row projection.  Negligible
transfers are then estimated with polynomial group and test counts.
The fixed-threshold argument following Equation~\eqref{cyl:tiny-energy}
discards tiny energies before the uniform exponent is applied.
A loss $s^{-\eps'}$ from a later outer limit is allowed only
when $\eps'/(j-i)\to0$; choose larger vanishing slack in the count
deletions.  These observations ensure that extraction produces genuine
polynomial subexperiments in the closure, not a formal limit object.
\end{proof}

\begin{proposition}[Bottom Planck floor]\label{cyl:planck-floor}
If the hybrid exponent satisfies $\Gamma_{\mathrm{hyb}}>\Gamma_{\mathrm{cl}}$
for the same weight, then every hybrid extremizing family at $\Gamma,p$
has
\begin{equation}\tag{Pl}\label{eq:Pl}
 H/(sf_1^2)=s^{o(1)}.
\end{equation}
At every fixed positive documented depth, also
$H/(\ell_i f_i^2)=s^{o(1)}$.
\end{proposition}

\begin{proof}
Admissibility gives the upper bound one.  Suppose the ratio is at most
$s^\xi$ for a fixed $\xi>0$.  Set $a=1-h$, with $h>0$ sufficiently
small relative to $\xi,\eta$, and $\delta=s^h$.  Apply the deletion and
chart partition in the proof of Proposition~\ref{cyl:matched-chain}
on $[a,1]$.  Each output has bounded chart overlap; summed chart energy
is at most $s^{-o(1)}E_a$, raw chart mass is at most
$s^{-o(1)}\kappa_aw^*$, and regularity down to the root Planck floor
is at most $s^{-o(1)}\kappa_aw^*$ in chart units.  The normalized output
normal width is $\delta^{D_*}$, whereas the final packet velocity width
is smaller by $s^{\xi/2}$.

Jitter as in the preceding proof.  The averaged initial row measure
can now be regarded as classical input regular down to width zero.
Above the hybrid floor use hybrid regularity.  Below that floor use
the density estimate $s^{-o(1)}\kappa_aw^*\delta^{-5l}$ and
Equation~\eqref{cyl:volume-bound}.  The fixed positive gap $\xi$ makes
the tiny-width gain dominate $\delta^{-5l}$ when $h$ is sufficiently
small.  All input cutoffs then hold classically for non-tiny chart
pieces.

We spell out why arbitrary hybrid masks inside the segment do not
invalidate the classical upper comparison.  In chart units use a
normal phase lattice of side $\zeta=\delta^{D_*+3}$.  At a final time
$t$, bin analyzed rows by $(N',D_t'-tN')$.  Along the fixed stack, all
packet drifts in these units are much smaller than $\zeta$ if $h$ is
small enough.  Thus, at every fixed exact index, contraction and
localization bound the output mass in one lattice cell by a constant
times initial row mass in finitely many neighboring cells, plus
negligible fiber energy.  Only polynomially many cells matter.

Reduce the assigned mass in each cell by its common excess factor so
that it fits this neighboring initial mass.  The reduction loses only
the negligible error after summing cells and integrating exact indices.
Divide once more by the bounded neighbor-overlap constant and allocate
the remaining assigned masses on those initial cells, separately for
each final time.  These are legitimate classical assignments: exact
base indices have the same endpoint, and a neighbor cell arrives within
$O(\zeta)$ in normal phase, less than the output positional tolerance
$\delta^{1+D_*}$.  Hence it lies in a bounded enlargement of the
original assigned test.  Keep the jitter parameter among the exact
indices throughout.  This is an energy comparison, not a positive
packet identity for a masked quantum evolution.

Applying the classical exponent chart by chart bounds the retained
hybrid cubic readout by
$s^{-o(1)}\delta^{1/2-\Gamma_{\mathrm{cl}}}E_a\sqrt{\kappa_a}$.
Saturation requires
$s^{o(1)}\delta^{1/2-\Gamma_{\mathrm{hyb}}}E_a\sqrt{\kappa_a}$,
a contradiction.  The prefix ending at any documented $i>0$ also
realizes $\Gamma,p$, by its read mass and count, so the same argument
applies there.
\end{proof}

\subsection{Curvature costs and shape profiles}

The matched gates control differences between successive shears. We now
ask how large the shear coefficients themselves must remain among
experiments with the extremal exponents $\Gamma,p$. Their least possible
growth rates will distinguish the configurations treated by the
lower-dimensional estimates from those requiring the entropy argument.

For a leaf label define
\begin{equation}\label{cyl:costs}
 R_2=\frac{r_1^2}{f_1}\bigl(|\beta(T_1)|+s|K|\bigr),\qquad
 R_1=\frac{r_1}{f_1}|\alpha+2\beta(T_1)x-2Kb|.
\end{equation}
The quantity $R_1$ measures the first-order horizontal variation of the
normal shear in units of $f_1$. The remaining variation has parabolic
order two: angular increments have size $r_1$, horizontal position
increments have size $sr_1$, and the tilted vertical increment has size
$sr_1^2$. Its size in normal velocity units is bounded by a constant
times $R_2$. The same ratios apply to the normal position condition,
whose width is $sf_1$.

Among families at $\Gamma,p$, minimize in the closure the exponent $c=c_2$
of an upper bound $s^{-c-o(1)}$ for $\max(1,R_2)$ on contributing
labels.  If $c=0$, also minimize the corresponding exponent $c_1$ for
$\max(1,R_1)$.  These are costs of the shear, not packet dimensions.
Section~\ref{sec:low} excludes simultaneous zero cost in the classical
model when $\Gamma>\eta/2$, and in the hybrid model when
$\Gamma>\Gamma_{\mathrm{cl}}$. The positive-cost alternatives are
reduced below to affine width profiles.

\begin{lemma}[Finite costs and compact profiles]\label{cyl:profiles}
The preceding minima have feasible finite values.  Normalize $r_1=f_1=1$.
Along successively finer matched meshes one may extract Lipschitz limits
\begin{equation}\label{cyl:profile-definition}
 u(i)=\lim\log_s r_i,\qquad n(i)=\lim\log_s f_i,\qquad
 b_t(i)=n(i)-tu(i)\quad(t=1,2).
\end{equation}
Here $u,n$ vanish at $1$ and $u$ is nonincreasing, hence nonnegative.
In the strict-gap hybrid case, $n(i)=(1-i)/2$.
Along any connecting chain from $i$ to $j$,
\begin{equation}\tag{A}\label{eq:A}
 \begin{split}
 |\beta_i(T_j)-\beta_j(T_j)|
        &\le s^{\min_{[i,j]}b_2-o(1)},\\
 |K_i-K_j|
        &\le s^{\min_{k\in[i,j]}(b_2(k)-k)-o(1)},\\
 |P_i(F)-P_j(F)|&\le s^{\min_{[i,j]}b_1-o(1)}.
 \end{split}
\end{equation}
One may also impose the individual curvature bound
\begin{equation}\tag{B}\label{eq:B}
 |\beta_i(T_i)|+\ell_i|K_i|
        \le s^{\min(-c,\min_{[i,1]}b_2)-o(1)}.
\end{equation}
\end{lemma}

\begin{proof}
Apply the one-step chart partition at the root and
Lemma~\ref{cyl:group-extraction}.  Subtraction of the full chart shear
bounds all reached leaf coefficients by fixed powers of $s^{-1}$ in
leaf units.  This proves finite feasibility for $c$.
When approaching finite $c$, the same chart partition can be used
without subtracting its quadratic terms.  After a base boost and
translation, the absolute quadratic coefficients are still bounded by
fixed powers in leaf units: conversion from $s|K|$ to full-duration $K$
costs at most $s^{-1}$.  Subtracting only the linear part then gives
finite feasibility for $c_1$ when needed.  Once $c_1$ is minimized, that
linear subtraction is unnecessary, since the leaf-center derivatives
and the bounded quadratic coefficients already bound the linear term.
Constant normal translations affect neither cost.  Base boosts and
translations preserve $R_2$ and the corresponding ray derivative $P$.
Each output knows boundedly many charts, so all these extractions retain
the relative count bound.

After $r_1=f_1=1$, Proposition~\ref{cyl:matched-chain} gives a uniform
Lipschitz bound for the interpolated width exponents and their endpoint
values.  Arzel\`a--Ascoli on finer meshes gives
Equation~\eqref{cyl:profile-definition}.  Equation~\eqref{eq:Pl}
at each positive depth, compared with its value at depth one, yields
$n(i)=(1-i)/2$ in the strict hybrid case, and continuity handles zero.
Telescoping Equation~\eqref{eq:M} proves Equation~\eqref{eq:A}.
For the first inequality evaluate each earlier affine coefficient at
the later time $T_j$, which lies within its interval; the corresponding
$K$ error is multiplied by at most that interval's length.  A finite
mesh contributes only $O(L_*\,\mathrm{mesh})+o(1)$ to the exponent.
Remove labels on no full geometric chain, whose transfer is negligible.

For Equation~\eqref{eq:B}, work first on a fixed fine mesh, with a
deletion tolerance larger than several mesh-error constants.  Suppose
the part of the $i$-gate violating the bound could preserve the terminal
exponent.  Ordinary-regularize just before that gate.  Its total raw
mass is at most $s^{-o(1)}E_i$.  At the next depth $j$, keep only labels
that can continue to a leaf in the graph imposed by the two curvature
comparisons in Equation~\eqref{eq:M}.  Their total read mass is at least
$s^{o(1)}E_j$.

For any predecessor on the high part, the leaf bound and telescoping
force its affine trajectory to satisfy
\[
 |\beta_i(T_1)|\le
       s^{\min(-c,\min_{[i,1]}b_2)-\mathrm{err}}
\]
at some terminal time $T_1$ within the reached $j$-interval.
The exponent $\mathrm{err}$ is smaller than the deletion tolerance.
An affine function whose value or duration-scaled slope violates
Equation~\eqref{eq:B} either never enters this range or enters it on
an interval much shorter than $\ell_j$.  Thus each high input row
predicts only $O(1)$ possible $T_j$ values.  For each $T_j$, restrict
the input to the rows predicting it and apply contraction; all other
transfers into the kept rows are negligible.  Summing gives total read
mass at most $s^{-o(1)}E_i$, contradicting $E_j/E_i=s^{-d(j-i)}$ and
$d>0$.  Therefore the low part preserves the output.  Make this deletion
at every mesh depth, then pass to finer meshes.  Strict tolerances and
the ordered closure limits justify the same conclusion for limiting
extremal configurations.
\end{proof}

\begin{proposition}[Shape constraints]\label{cyl:shape-constraints}
If $c>0$, then
\[
 b_2(i)\ge c(i-1),
\]
with equality at every depth when $c\ne1$.
If $c=0<c_1$, then $b_1(i)=c_1(i-1)$, and the quadratic terms of every
gate shear may be replaced by their linearization at the test center,
with vanishing exponent enlargement of its normal widths.
\end{proposition}

\begin{proof}
Write $b=b_2$.  At a positive depth $i$ where $b(i)=\min_{[i,1]}b$,
the prefix realizes $\Gamma,p$ by read mass and count.  Multiplying
Equation~\eqref{eq:B} by $r_i^2/f_i=s^{-b(i)+o(1)}$, its normalized
terminal curvature cost is at most $(b(i)+c)_+/i$.  Minimality of $c$
therefore implies
\[
 ci\le(b(i)+c)_+.
\]
Since $c>0$ and $i>0$, this yields $b(i)\ge c(i-1)$ at every positive
future minimum.  For any other depth, a later future minimum gives the
same lower bound, because $c(i-1)$ increases with $i$; continuity handles
zero.

Comparison with root labels using Equations~\eqref{eq:A} and
\eqref{eq:B} gives
\begin{equation}\label{cyl:K-profile}
 |K_i|\le s^{\min(-c,\min_{k\in[0,i]}(b(k)-k))-o(1)}.
\end{equation}
Suppose first $0<c<1$ and at some interior $i$ the shape inequality is
strict.  The lower comparison line for $b(k)-k$ is strictly decreasing.
By continuity, Equation~\eqref{cyl:K-profile} improves
$\ell_i|K_i|\le s^{c(i-1)-o(1)}$ by a fixed power at this $i$.
The minimum of $b$ on $[i,1]$ likewise strictly exceeds $c(i-1)$.
Consequently, for some $\eps>0$, every connected leaf satisfies
\[
 |\beta_i(T_i)-\beta_1(T_1)|\le s^{c(i-1)+\eps},
 \qquad s|K_1|\le s^{-o(1)}.
\]
At the regularized $i$-gate, bin $\beta_i(T_i)$ at that precision and
subtract the rounded constant quadratic shear.  Every reached leaf
knows only boundedly many bins.  Group extraction gives an experiment
on $[i,1]$ at $\Gamma,p$ whose terminal curvature cost, after dividing
exponents by $1-i$, is strictly smaller than $c$.  This contradicts its
minimality.

For $c>1$, a strict shape inequality at $i$ instead gives
\[
 \min_{k\in[i,1]}(b(k)-k)>c(i-1)-i,
\]
since its lower comparison line is strictly increasing.  Equation~\eqref{eq:A}
then compares the pair $(\beta(T_i),K)$ for the $i$-label and a leaf to
precisions $s^{c(i-1)+\eps}$ and $s^{c(i-1)-i+\eps}$, respectively;
extrapolating $\beta$ to $T_i$ costs at most $\ell_i$ times the $K$
error.  Bin this jointly predictable pair and subtract its polynomial
trajectory.  Group extraction again lowers $c$, a contradiction.

Now let $c=0<c_1$.  A future minimum of $b_1$ is also a future minimum
of $b_2=b_1-u$, since $u$ is nonincreasing.  At such a depth,
Equation~\eqref{eq:B} makes the normalized prefix curvature cost zero.
On a connecting chain the center derivatives at depths $i$ and $1$
differ by at most $s^{\underline b_{1,i}-o(1)}$, where
$\underline b_{1,i}=\min_{[i,1]}b_1$.  To see this, first compare them
at the actual ray by Equation~\eqref{eq:A}.  Moving to the center at
$i$ costs at most the curvature bound times $r_i$; its exponent is
at least $\underline b_{1,i}$ because
\[
 \min(0,\min_{[i,1]}b_2)+u(i)\ge\underline b_{1,i}.
\]
The leaf-center change has zero cost.  Lexicographic minimality at
positive future minima therefore gives
$b_1(i)\ge c_1(i-1)$, as in the curvature argument.  At any strict
interior point, bin and subtract the jointly predictable linear
coefficient.  The extracted tail keeps zero curvature cost while
strictly lowering $c_1$.  Thus equality holds everywhere.

Both $b_1$ and $b_2=b_1-u$ are now nondecreasing.  Equation~\eqref{eq:B}
therefore bounds curvature at depth $i$ by $s^{b_2(i)-o(1)}$.
Taylor expansion at the base center leaves the coefficient $P_i(x,b)$
times the increments in $X$ and $B_{T_i}$ in the respective normal
conditions.  Every remaining term is quadratic in horizontal increments
or linear in the tilted vertical increment, so has size at most
$s^{-o(1)}f_i$ in velocity and $s^{-o(1)}\ell_i f_i$ in position.
The existing row supports can consequently be kept.
\end{proof}

\subsection{Affine reduction}

The preceding constraints need not make both width profiles affine.
We obtain affine profiles by restricting to a short depth interval near
a common differentiability point and normalizing that interval. The
extraction must preserve the extremal exponents and minimal shear costs;
otherwise the new profiles would not describe the configurations still
to be excluded.

\begin{proposition}[Affine extremal configurations]\label{cyl:affine}
If at least one relevant minimal cost is positive, one can pass to
extremizing documented configurations with
\begin{equation}\tag{F}\label{eq:F}
 u(i)=J(1-i),\quad J\ge0,\qquad
 \begin{cases}
 b_2(i)=q(i-1),&c>0,\\
 b_1(i)=c_1(i-1),&c=0<c_1.
 \end{cases}
\end{equation}
For $c\ne1$, $q=c$.  For $c=1$, either $q=1$ or $q<1$; no
nonnegative lower bound on $q$ is assumed.  In the latter option,
\[
 |\beta_1(T_1)|\le s^{-1-o(1)},\qquad
 s|K_l|\le s^{-o(1)}\quad\hbox{at every gate }l.
\]
For $0<c\le1$ and $q=c$, one has
$\ell_l|K_l|\le s^{b_2(l)-o(1)}$, and the own-time values
$\beta_i(T_i),\beta_l(T_l)$ on forward connections differ by at most
$s^{b_2(i)-o(1)}$.  For $c>1$, compare instead
$\beta_i(T_i),K_i$ with $\beta_l(T_i),K_l$ at respective precisions
$s^{b_2(i)-o(1)}$, $s^{b_2(i)-i-o(1)}$.
The linear alternative retains its individual curvature bounds and the
center-derivative comparison at precision $s^{b_1(i)-o(1)}$.

All relevant terminal minimal costs, gate counts, read-mass bounds,
Equations~\eqref{eq:M} and~\eqref{eq:A}, and the permission to insert
ordinary regularization persist.  Base support, significant normal
support, and reachable shear parameters are bounded by fixed powers
in every normalized inner experiment.
\end{proposition}

\begin{proof}
Choose an interior common differentiability point of the Lipschitz
profiles and set $J=-u'$ and $q=b_2'$ when relevant.  If $c=1$ and the
shape is not the equality line, there are points with $q<1$: otherwise
absolute continuity and the endpoint value would contradict the strict
gap somewhere above $i-1$.  If the shape is the equality line, take
$q=1$.  Take shrinking segments $[A,B]$ about this point, put
$h=B-A$ and $\Delta=s^h$, and normalize duration and terminal widths
to $(\ell_A,r_B,f_B)$.  For each fixed $h$, take the mesh fine enough
and all preceding closure errors small enough that their exponents are
$o(h)$.  Differentiability then gives Equation~\eqref{eq:F} uniformly
in the normalized depth.  It remains to justify extracting actual
experiments without losing the costs or compatibility.

At the ordinary-regularized $A$-gate partition rows by a containing
lattice base chart in $r_B,\ell_A$ units, and boost and translate to
its center.  Every later connected label at depth $l\le B$ knows only
$s^{-o(h)}$ possible charts, since $r_l\le s^{-o(h)}r_B$ and its base
backflow is bounded in those units.  Use the following additional bins,
written temporarily in the old physical width exponents:
\begin{enumerate}[label=\textup{(\roman*)}]
\item If $c>1$, bin $(\beta_A(T_A),K_A)$ at widths
$s^{b_2(A)}$, $s^{b_2(A)-A}$, and subtract the rounded polynomial shear.
Equation~\eqref{eq:A}, extrapolated to $T_A$, makes these bins known
to every connected later label up to $s^{-o(h)}$ choices.
\item If $0<c<1$, bin only $\beta_A(T_A)$ at width $s^{b_2(A)}$ and
subtract that constant quadratic coefficient.  The bounds for $K$ in
Proposition~\ref{cyl:shape-constraints} make every later own-time
coefficient predict the bin.
\item If $c=1$ and $q\le1$, bin the pair at widths
$s^{b_2(B)+A-B}$ and $s^{b_2(B)-B}$, and subtract the rounded trajectory.
On this segment $b_2(k)-k$ is affine nonincreasing up to $o(h)$, so
Equation~\eqref{eq:A} again gives the required prediction.
\item In the linear alternative subtract no curvature.  Bin the
$A$-label's center derivative at width $s^{b_1(A)}$ and subtract that
linear shear; Proposition~\ref{cyl:shape-constraints} supplies the
joint prediction from every connected later label.
\end{enumerate}
The subtractions are made after centering and include free constant
terms.  All residual curvature data have bounded-power size in $\Delta$
units.  For example, in case (iii), the uncertainty of $K$ contributes
at most $s^{b_2(B)+A-B-o(h)}$ to the terminal quadratic coefficient,
and at most $s^{b_2(B)-o(h)}$ to its terminal duration-scaled $K$ term.
This gives terminal cost at most one, and when $q<1$ yields precisely
$s|K_l|\le s^{-o(1)}$ after normalization.  Cases (i), (ii), and (iv)
give the other displayed bounds directly.  Their terminal cost is at
most its prescribed minimum; minimality prevents a strict improvement.

When $c>0$, also bin and subtract a linear shear at a sufficiently
large-power tolerance in $\Delta$ units.  This has no effect on the
curvature costs.  Indeed compare derivatives first on the connecting
ray by Equation~\eqref{eq:A}; $\min b_1$ differs from $b_1(B)$ by
$O(h)$.  The already bounded residual curvature permits recentering
at the common base chart, still at a fixed-power cost.  Every later
label therefore predicts the additional bin up to $s^{-o(h)}$ choices.
In the linear alternative these derivatives are already bounded.
Finally bin the actual residual normal phase at $A$ at a sufficiently
large-power tolerance and remove its constant velocity and position.
Base backflow, bounded residual coefficients, and bounded packet drifts
make each later reached label predict these bins with the same small
multiplicity.  The starting labels themselves also meet only that many
groups.

Apply Lemma~\ref{cyl:group-extraction}, including the relative count
conditions at the finitely many intermediate depths.  Transform the
selected experiment by the indicated symmetries and remove labels
unreachable from the selected group.  Negligible errors are evaluated
on each inner polynomial setup before its outer limits.  Gate counts
are now relative to the extracted experiment's own energy and inherited
regularity.  Its raw starting mass is at most that regularity times a
fixed power of $\Delta^{-1}$, by its bounded support; synthesized tails
are negligible.  Its energy is at least that regularity times an inverse
fixed power for a near-extremal extracted group: otherwise contraction
onto the normalized contributing output tests would make its normalized
readout negligible by Equation~\eqref{cyl:tiny-energy}.  Thus all cutoff
requirements hold after normalization.

If a documented root gate is needed, restore it on the extracted
post-state.  Classically keep the old row ownership.  In the hybrid case
reanalyze and retain rows having a nonzero original group row, in the
same exact fiber, within $s^{-\eps}$ packet widths.  Use essential
neighborhoods, defined by positive fiber mass in the relevant box, and
assign one witnessing original label.  Gaussian localization makes
projection and synthesis approximate the identity on this input up to
negligible error.  Enlarge the root normal widths accordingly.  Choose
$\eps=o(h)$ smaller than the truncation slack at the first mesh step.
Every connection from the restored gate then extends geometrically to
an original supporting row, so the original compatibility estimates
remain valid with vanishing normalized errors.

The mesh in the normalized depth may now tend to zero, always after
the inner scale limits.  In relative count selections take vanishing
$\Delta$-exponent slack larger than the total overlap and extremality
errors.  This proves realization of $\Gamma,p$ and the same terminal
cost minima in the closure.  All arguments concerned finite masks and
finite meshes; no infinite product of operators has been taken.
\end{proof}

\begin{remark}[Forward preparation of later observations]\label{cyl:forward-preparation}
One may ordinary-regularize all available gates on each finite affine
mesh in forward order, interleaving any prescribed finite row statistics
at that gate, before flattening the final readout.  For any subsequently
chosen fixed interior depth, a nearest available gate then already has
the benchmarks in Equation~\eqref{cyl:benchmarks}.  Its prefix and tail
have positive limiting depth gaps, so
Proposition~\ref{cyl:split-saturation} applies.  This procedure neither
asserts relative saturation on every shrinking adjacent mesh gap nor
transfers profiles retroactively through new earlier masks.
\end{remark}

\subsection{Causal laws of the raw rows}

Fix finitely many observed gates, with ordinary pre-gate regularization
prepared in forward order. Let $\Prob^i$ be the probability measure
on their retained pre-synthesis rows at depth $i$, including time, exact
indices, and documented label, obtained by dividing raw row mass by its
total. At depth one use the final assigned mass.
Proposition~\ref{cyl:matched-chain} and pre-gate energy give
\begin{equation}\label{cyl:law-normalization}
 s^{o(1)}E_i\le\text{normalizing mass}\le s^{-o(1)}E_i.
\end{equation}
Nearest mesh gates may represent requested depths, with vanishing
exponent ambiguity. All these laws belong to actual finite inner setups.

Correspondence predecessors below are taken in the current starting-state
essential support and through the pre-existing intermediate physical
masks. We do not require membership in the original state's intermediate
supports: restricting the input may undo cancellation at those cuts.

We next compare the mass of an event of later rows with the mass of its
possible predecessors. If all predecessors of a later event $U$ belong
to an earlier row set $S$, locality permits us to synthesize only $S$.
The cylinder exponent then bounds the output on $U$. This is the rarity
estimate below; it does not require a hybrid output sample to have a
sampled ancestral packet.

\begin{lemma}[Rarity transport]\label{cyl:rarity}
Let $i<j$. Let $U$ be a measurable event of retained $j$-rows and
$S$ a measurable event of retained $i$-rows. Suppose every correspondence
predecessor in the active essential support of a row in $U$ belongs
to $S$, uniformly at the selected truncation. Then
\begin{equation}\tag{Sp}\label{eq:Sp}
 \Prob^j(U)\le s^{-o(1)}\Prob^i(S)^{2/3}+\mathrm{negl}.
\end{equation}
\end{lemma}

\begin{proof}
Restrict synthesis at $i$ to $S$. Its summed input energy is at most
$s^{-o(1)}\Prob^i(S)E_i$, and its regularity remains at most
$s^{-o(1)}\kappa_i$ in each interval. By locality, its output on $U$
agrees with the original output up to negligible error: input outside
the active essential support has zero state. The tail
exponent bounds its cubic readout by
\[
 s^{(1/2-\Gamma)(j-i)-o(1)}
                 \sqrt{\kappa_i}\,\Prob^i(S)E_i.
\]
H\"older and the $j$-gate weight count give
\[
 \begin{split}
 \sum_{\lambda\text{ on }U}m_\lambda
 &\le
 \left(\sum_\lambda m_\lambda^{3/2}w_\lambda^{-1/2}\right)^{2/3}
 \left(\sum_\lambda w_\lambda\right)^{1/3}\\
 &\le s^{-o(1)}E_j\Prob^i(S)^{2/3}.
 \end{split}
\]
The scale powers cancel by $2d=2\Gamma+p-1$.
Divide by Equation~\eqref{cyl:law-normalization}.
Superpolynomially small input pieces are treated by
Equation~\eqref{cyl:tiny-energy}; polynomial weights and counts absorb
the negligible operator errors. At final readout the restriction to
$U$ means partial assignment, so the same argument applies.
\end{proof}

For the classical model a final sample of nonzero mass does have an
exact-particle path through earlier pointwise masks, with ownership at
the ancestor disjoint gates. We use this additional observation only
in the classical case. Lemma~\ref{cyl:rarity} is the substitute that
is valid in both models.

To apply this estimate when a row is marked by the existence of a
successful completing path, we need measurable versions of the marked
rows and their predecessor sets. Use a Borel inner version of the end
event and a Borel outer version of all its geometric predecessors.
The following lemma preserves the completion of every retained end row
and the pointwise predecessor inclusion required by locality. Its proof
is given at the end of this subsection.

\begin{lemma}[Measurable completion events and causal restrictions]
\label{cyl:borel-paths}
Fix an inner scale, a finite stack of cuts, a finite query list, and a
correspondence truncation.
\begin{enumerate}[label=\textup{(\roman*)}]
\item The side and row data admit standard Borel codes in which the
primitive correspondences, time ownership, physical masks, analyzed-state
essential supports, and assignment supports have Borel versions. All
energetically active paths occur in a Borel finite-path relation. Its coordinate
projections, including completion events, are analytic and hence
measurable in the completion of every Borel row law
\cite[Section~3, Theorem~2]{BertsekasShreveBorel}.

\item For an analytic row event $A$ and a finite Borel row law $\mu$,
there are Borel versions $A^-\subset A\subset A^+$ with
$\mu(A^+\setminus A^-)=0$. For a pre-synthesis mask, $\mu$ is the
\emph{full analyzed row-energy law} $\|V_i g_i\|^2$ in the active
intervals, before fractional assignments. An assigned law suffices only
for a terminal readout. An inner version of a completion event retains
a completion for every one of its rows.

\item Let $B$ be a Borel inner version of a marked end-row event,
intersected with the current end-state positive-norm set $H_j$; it has
the same full analyzed end-row mass as the marked event. Let $S$ consist of all predecessors
of $B$ under the unrestricted truncated geometric correspondence from
the starting gate to the marked gate. This relation uses the current
starting-state essential support and the pre-existing intermediate
physical masks. It does not impose the success event that selected $B$
again, or the original state's intermediate supports. The set $S$ is
analytic and has Borel versions
\[
 S^-\subset S\subset S^+,
 \qquad \mu(S^+\setminus S^-)=0
\]
under the full analyzed starting-row energy law. Thus $S^+$ contains
every geometric predecessor pointwise. Synthesis restricted to $S^+$ and
$S^-$ gives identical states, and the two restrictions have equal
analyzed input energy.
This assertion concerns the physical state; a completing witness may
have zero conditional probability under a separate path law.

\item With finitely many possible witness bins, the marked end rows
may first be disjointized by bin and given Borel inner versions in
$H_j$. Each retained row still has a completion in its chosen bin.
Form each bin's predecessor set by the unrestricted correspondence in
\textup{(iii)}. If path concatenation places this set in a Borel allowed
starting-bin set $P_b$, both predecessor versions may be chosen inside
$P_b$. The original starting-bin partition and its overlap count are
then preserved. Use the outer versions for pointwise locality and the
inner versions for physical restrictions and input-energy accounting.

\item The versions may be chosen jointly in side and row coordinates,
and simultaneously for any fixed finite or countable family of full
analyzed and comparison laws. After a new restriction or a row law not
dominated by those laws is formed, choose versions under that law before
the next projection. Complements of analytic events are used only as
completed-measurable events, or given Borel inner versions before
entering another path relation. Conditional assertions obtained by
disintegration hold for almost every side value; a uniform assertion
requires a separate finite-grid or quantitative exceptional-side
argument.
\end{enumerate}
\end{lemma}


\begin{lemma}[Forward flattening of scores]\label{cyl:flattening}
At finitely many gates in forward order, prescribe finitely many
finite-valued measurable row names, each having polynomially many
values.  Conditional names may also condition on a common arbitrary
exact field $y$.  One may retain subpower choices of row restrictions,
preserving the terminal exponent, so that the prescribed log probability
scores have deterministic limiting values.  Their conditional versions
are evaluated under the actually retained law.  Further restrictions
of that same law of probability $s^{o(1)}$ preserve these deterministic
values, outside events of power-small probability at every fixed
positive score tolerance.
\end{lemma}

\begin{proof}
Compute the probabilities under the gate's preselection law, and bin
their nonnegative $\log_s$ scores.  Scores beyond a sufficiently large
power occur on power-small mass, by the finite-name capacity bound,
also after conditioning on $y$.  Their contribution cannot preserve
the terminal exponent: restrict to those rows and use the tail exponent,
or Lemma~\ref{cyl:rarity}.  There remain subpower many bins for each
fixed finite request.  The coherent triangle inequality selects a
class preserving the output.  Such a class has probability $s^{o(1)}$,
since a power-small fraction would again lose the output exponent.
At final assignment one uses
$(\sum_{b=1}^k m_b)^{3/2}\le k^{1/2}\sum_{b=1}^k m_b^{3/2}$.

Here is the required change-of-law calculation.  If $Q=P(\,\cdot\mid A)$
with $P(A)=s^{o(1)}$, then for any measurable name or conditioning field
$F$,
\[
 r_F=\frac{dQ_F}{dP_F}\le P(A)^{-1},\qquad
 Q\{r_F<s^\eps\}\le s^\eps.
\]
The first inequality follows from $Q(B)\le P(B)/P(A)$; the second
follows by integrating $r_F$ over the indicated event. Write $D$ for one
prescribed name and $C$ for its finite-valued conditioning data. Apply
these inequalities to $(y,C,D)$ and $(y,C)$. Their ratio is the likelihood
ratio for the conditional probability of $D$ given $(y,C)$.
Its $\log_s$ is arbitrarily small outside a power-small event.
Thus the scores selected before conditioning have the same limiting
values under the retained law.  This is also
Lemma~\ref{lem:fw-restriction}.

Make later gate choices only after earlier ones have been realized.
Countably many eventual requests mean increasing finite lists followed
by subsequences; they do not mean infinitely many projections on one
state.  Dense-query extension is legitimate only when the corresponding
names have explicitly verified small-list relations as the query
parameters approach one another.
\end{proof}

\begin{lemma}[Transfer of conditional profiles]\label{cyl:profile-transfer}
Suppose names on the $i$- and $j$-row laws have the following finite-list
relations on every connecting path.  Given a common exact variable $y$,
each conditioning name $C_i$ determines at most $s^{-o(1)}$ possibilities
for $C_j$, and conversely.  Given $y$ and a related pair of conditioning
names, the observed names $D_i,D_j$ likewise have bidirectional
$s^{-o(1)}$ lists.  If
\[
 \log_s\Prob^i(D_i\mid y,C_i)=z+o(1)
\]
outside power-small events at every fixed tolerance, then the same
statement holds with $i$ replaced by $j$.  All lists must be measurable
and uniformly valid for the correspondence geometry.
\end{lemma}

\begin{proof}
Fix $\eps>0$.  At each $(y,C_j)$ the set of $D_j$ atoms of probability
greater than $s^{z-\eps}$ has cardinality at most $s^{-z+\eps}$.
Pulling this list back through the conditioning and observed-name lists
gives at most $s^{-z+\eps-o(1)}$ possible $D_i$ values per $(y,C_i)$.
Outside the exceptional predecessor set, each such value has conditional
probability at most $s^{z-\eps/4}$.  Their union therefore has
power-small $\Prob^i$ mass.  Lemma~\ref{cyl:rarity} makes the
corresponding large-atom event power-small at $j$.

For the other direction, the good $D_i$ atoms of conditional probability
at least $s^{z+\eps/4}$ have at most $s^{-z-\eps/4}$ values per
$(y,C_i)$.  Forward their observed-name lists and use the backward
conditioning lists.  The resulting set has at most
$s^{-z-\eps/4-o(1)}$ $D_j$ values per $(y,C_j)$.  Missing this set
requires an exceptional predecessor and is power-small by
Lemma~\ref{cyl:rarity}.  Within the set, atoms below $s^{z+\eps}$
have total conditional probability at most $s^{3\eps/4-o(1)}$ by
counting.  This proves the claim.  The same proof within one law shows
invariance under changes of discretization with the stated list bounds.
\end{proof}

\begin{lemma}[Time prediction]\label{cyl:time-prediction}
Let $i<j\le l$.  Suppose a conditioning name at $l$ has relay names
at $j$ and $i$ such that an $i$-row gives at most $s^{-o(1)}$ possibilities
for the relay at $j$, and that relay gives at most $s^{-o(1)}$
possibilities for the target conditioning name at $l$.  The relations
must hold on all connecting paths; an additional common exact variable
$y$ is permitted.  Then any measurable list of at most
$s^{-d(j-i)+\eps}$ candidates for $T_j$ per target conditioning datum
has power-small probability under the $l$-law, for fixed $\eps>0$.
Consequently the conditional log probability score of the time choice
is at least $d(j-i)-o(1)$, outside power-small events at every fixed
tolerance.
\end{lemma}

\begin{proof}
Mark $j$-rows whose relay predicts a target datum whose proposed time
list contains their $T_j$.  An $i$-row can reach marked rows at no more
than
$N=s^{-d(j-i)+\eps-o(1)}$ possible times.  For each such time restrict
the input to its possible predecessor rows, and use contraction and
locality.  Summing squared norms over times costs at most $N$ per input
row, so marked mass at $j$ is at most
\[
 s^{-o(1)}NE_i=E_j s^{\eps-o(1)}.
\]
If $l>j$, apply Lemma~\ref{cyl:rarity} to transport this exceptional
event; the final probability is at most $s^{2\eps/3-o(1)}$.
For $l=j$ divide directly by Equation~\eqref{cyl:law-normalization}.
To obtain the conditional score claim, take the proposed list to be
the time atoms larger than $s^{d(j-i)-\eps}$; there are at most the
stated number.  Finite grids and removal of the resulting exceptional
events give the corresponding prefix mass bounds.  Availability of the
relay at the tested prefix is essential to the input multiplicity
estimate.
\end{proof}

\begin{lemma}[Recoverable phase names]\label{cyl:phase-name}
Let a finite phase-cell name be observed on a later law.  Suppose that,
for each name, all its predecessors at an ordinary-regularized gate
$i$ are covered by $s^{-o(1)}$ time/test pairs of weights at most
$s^{-o(1)}w$.  Its typical unconditional log probability score is at least
\[
 \log_{1/s}\frac{E_i}{\kappa_iw}-o(1).
\]
If also $w=w_i s^{o(1)}$ and every documented $i$-label predicts at most
$s^{-o(1)}$ names on all reachable rows, this lower bound is an equality
in the limit.
\end{lemma}

\begin{proof}
Write $z=\log_{1/s}(E_i/(\kappa_iw))$.  A list of at most
$s^{-z+\eps}$ names has predecessor mass at most
$s^{\eps-o(1)}E_i$ by regularity and the covering hypothesis.
Lemma~\ref{cyl:rarity} makes its probability power-small on the later
law.  Apply this to the list of atoms larger than $s^{z-\eps}$ to
obtain the lower score bound.

For the reverse bound, the $i$-gate weight count and the prediction
hypothesis cover all reachable later rows by at most
$s^{-o(1)}E_i/(\kappa_iw_i)=s^{-z-o(1)}$ names.  Any remaining rows
have negligible probability by locality and Lemma~\ref{cyl:rarity}.
Within a list of that size, atoms smaller than $s^{z+\eps}$ have total
mass at most $s^{\eps-o(1)}$.  This proves equality.
\end{proof}

The last three lemmas require actual geometric recoverability and finite
measurable lists.  Counting formal coordinates of a name, without
proving these relations, does not establish their hypotheses.

\paragraph{Measurability of correspondence events.}
We finish by proving the measurable-version statement used in the causal
arguments above.

\begin{proof}[Proof of Lemma~\ref{cyl:borel-paths}]
Only finitely many cuts and queries occur at this inner setup.  Keep
the exact-index spaces themselves in the row code, together with the
Euclidean phase, time, and finite label coordinates; a standard Borel
space admits an injective Borel realization in a Polish space.  Thus
equality of a preserved exact index is a Borel diagonal condition and
distinct indices are never identified.  Use the measurable
separable-Hilbert coordinates of the direct integral to represent
the finitely many operator masks.  Their countably many matrix
coefficients have Borel versions under the reference measure; on the
common exceptional null set set the mask to zero.  This leaves the
pointwise contraction and the multiplication operator unchanged.
At every fixed cut and interval, take a Borel measurable-field
representative of the analyzed state and let $H_i$ be its Borel
positive-norm set. The full analyzed row-energy law is concentrated
on $H_i$. Changing the representative on a reference-null set, or
inserting $\mathbf1_{H_i}$ before synthesis, leaves the physical state
unchanged. In an $i$-to-$j$ predecessor test use the current input
state's $H_i$ and the marked end event as vertex conditions. At
intermediate cuts use the Borel masks and geometric correspondences,
not the original state's support: a restriction at $i$ can undo
cancellation there. If a support condition for a particular
propagated component is needed, recompute its Borel representative
after that physical restriction. Thus ambient zero-state starting
rows are not witnesses, without omitting possible intermediate
transfer from the actual restricted input.
Fractional assignment densities relative to the actual row measures
likewise have Borel versions and give the same assigned measures.
Finite time ownership, test inequalities, exact-index equality, and
the truncated Gaussian correspondences in \eqref{cyl:correspondence}
are Borel in the code.
Their finite intersection in the full path space is therefore Borel.

Projection of this Borel relation is analytic.  Analytic and
coanalytic sets are universally measurable, so completion of the
active Borel law gives Borel inner and outer versions of equal mass.
For finitely or countably many laws use a weighted sum of their
normalizations to choose common versions.  In particular, taking the
inner version of the marked end rows preserves each retained row's
literal witness. Intersect this inner version with the current
end-state $H_j$, which is Borel and has full end-row energy. Taking
the outer version of their predecessor set
preserves inclusion of \emph{all} geometric predecessors.  One may
also choose a Borel inner predecessor version $S^-\subset S$ of equal
full analyzed input mass.  The difference $S^+\setminus S^-$ carries
zero current pre-synthesis input state, so synthesis restricted to
$S^+$ or $S^-$ is identical.
This last equality concerns the actual state, not an arbitrary
probability law on completing paths; a witness may have zero
conditional mass.

When finitely many witness bins are used, first form the analytic
possible-end-row set for each bin.  Lexicographically disjointize
them, take Borel inner versions under the full analyzed end law,
intersect each with the current end-state $H_j$, and then
project the unrestricted starting-to-end geometric correspondence
for each disjoint marked group. Take Borel inner and outer
predecessor versions $S_b^-\subset S_b\subset S_b^+$ of equal full
analyzed starting-row energy. If path concatenation places $S_b$
in a Borel allowed starting-bin set $P_b$, intersect both versions
with $P_b$; the
original starting-bin partition and its overlap count are unchanged.
Use $S_b^+$ for pointwise locality and $S_b^-$ for the physical
state and input-energy accounting. Every kept end row still has a
witness in its assigned bin.  A later
physical restriction uses its Borel version in a newly formed
finite-path relation.  If an approximate inner version is convenient,
losing row energy at most $\xi^2$ changes the restricted input state
by at most $\xi$ in norm; the following contractions preserve that
bound.  Choose $\xi$ smaller than every required power of $s$ before
the exponent limits, so polynomially many descendants and tests add
only negligible errors.  Finally, choose joint versions under the
side-row law and disintegrate.  This proves almost-everywhere, not
pointwise, conditional assertions.
\end{proof}
